\documentclass[11pt,a4paper]{article}
\usepackage[T1]{fontenc}
\usepackage[utf8]{inputenc}
\usepackage{lmodern,amsmath,amssymb,amsthm,mathtools}
\usepackage[a4paper,margin=26mm,headheight=14pt]{geometry}
\usepackage{microtype,booktabs,array,longtable,enumitem,xcolor}
\usepackage{fancyhdr,hyperref}
\definecolor{ink}{HTML}{183447}
\definecolor{accent}{HTML}{226D7A}
\hypersetup{colorlinks=true,linkcolor=ink,citecolor=accent,urlcolor=accent,
 pdftitle={Order Derivatives Recover Hidden Polylogarithmic Transseries},
 pdfauthor={Research report prepared for Vladimir Reshetnikov}}
\pagestyle{fancy}\fancyhf{}
\fancyhead[L]{\small Polylogarithmic order jets and hidden transseries}
\fancyhead[R]{\small 12 October 2026}\fancyfoot[C]{\thepage}
\setlength{\parskip}{4pt}\setlength{\parindent}{1.15em}
\setlist{itemsep=3pt,topsep=4pt}
\numberwithin{equation}{section}
\newtheorem{theorem}{Theorem}[section]
\newtheorem{proposition}[theorem]{Proposition}
\newtheorem{lemma}[theorem]{Lemma}
\newtheorem{corollary}[theorem]{Corollary}
\theoremstyle{definition}\newtheorem{definition}[theorem]{Definition}
\newtheorem{example}[theorem]{Example}
\theoremstyle{remark}\newtheorem{remark}[theorem]{Remark}
\newcommand{\C}{\mathbb C}\newcommand{\R}{\mathbb R}
\newcommand{\N}{\mathbb N}\newcommand{\Z}{\mathbb Z}
\newcommand{\ii}{\mathrm i}\newcommand{\e}{\mathrm e}
\newcommand{\dd}{\,\mathrm d}\newcommand{\Li}{\operatorname{Li}}
\newcommand{\Res}{\operatorname{Res}}\newcommand{\PV}{\operatorname{PV}}
\newcommand{\FP}{\operatorname{Fp}}\newcommand{\Log}{\operatorname{Log}}
\newcommand{\B}{\mathcal B}\newcommand{\Sp}{\mathcal S_+}
\newcommand{\Sm}{\mathcal S_-}\newcommand{\Med}{\mathcal M}
\newcommand{\eps}{\varepsilon}\newcommand{\muh}{\widehat\mu}
\newcommand{\abs}[1]{\left|#1\right|}
\newcommand{\coef}[2]{[#1^{#2}]}
\newcommand{\pin}{\texttt{7bd1ce58abe522b83cd3368bbafa1be97837b75d}}
\begin{document}
\hypersetup{pageanchor=false}
\begin{titlepage}
\setlength{\parindent}{0pt}
\vspace*{9mm}
{\color{accent}\large SPECIAL FUNCTIONS AND TRANSSERIES\hfill RESEARCH REPORT}\par
\vspace{19mm}
{\color{ink}\fontsize{29}{35}\selectfont\bfseries
Order Derivatives Recover\\[3pt]
Hidden Polylogarithmic\\[3pt]
Transseries}\par
\vspace{10mm}
{\Large Exact completion, resonance jets,\\[3pt]
and finite reflection rigidity}\par
\vspace{12mm}
{\large Prepared for Vladimir Reshetnikov and the ProveIt research program}\par
\vspace{4mm}
{12 October 2026}\par
\vspace{12mm}
\begin{abstract}
We study shifted Lambert--polylogarithm sums
$F_s(t,a)=\sum_{\ell\ge0}\Li_s(\e^{-(\ell+a)t})$.
At each positive integer order $m$, a reflection combination has a finite
algebraic expansion but a generally nonzero, exactly computable
exponential remainder. Its first order derivative restores a divergent
series with double Borel poles. With an explicit Borel normalization, the
Stokes discontinuity of this derivative is exactly $\pi\ii$ times the
hidden remainder. We also identify the derivative itself: its linear
median requires a convergent dual correction with logarithmic divisor
weights. This identification is proved by Mellin inversion, not inferred
from the discontinuity.

For arbitrary finite sets of real shifts, we classify all terminating
integer-order tails and give finite Fourier--M\"obius certificates.
A nonzero reflection-compatible measure on $r$ atoms has a hidden action
among the first $r$ lattice actions; midpoint grids show the bound is
sharp. More generally, the joint kernel of the zeroth and first order
 tails, on finite measures, is precisely the span of Lebesgue measure.
An all-order collision formula shows how reflection removes the explicit
Stieltjes term and lowers the logarithmic degree at every order jet.
Proofs, a source audit, and reproducible exact and numerical checks are
included. Scalar Lambert resurgence and the order-one $q$-product
completion are credited to earlier literature and repository work.
\end{abstract}
\vfill
\noindent\small\textbf{Status.} Conventional mathematical proofs are supplied.
The results have not been independently refereed or proof-assistant
formalized. The intended contribution is the combined all-order,
all-real-shift formulation and its finite rigidity consequences;
first-priority in the literature is not asserted.
\end{titlepage}
\hypersetup{pageanchor=true}\setcounter{page}{1}
\begingroup
\small\setlength{\parskip}{0pt}
\tableofcontents
\endgroup
\newpage

\section{The problem and the contribution}
\label{sec:scope}

A finite asymptotic expansion need not determine an analytic function.
In the family studied here, reflection can annihilate every coefficient
of a divergent tail while leaving infinitely many exponentially small
terms. The point of this report is to recover those terms from a
\emph{specified analytic deformation}: differentiation in the
polylogarithm order. We prove both the recovery formula and the analytic
completion of the differentiated function.

The variables and the principal observables are
\begin{align}
 F_s(t,a)&=\sum_{\ell\ge0}\Li_s(\e^{-(\ell+a)t})
       =\sum_{k\ge1}\frac{\e^{-akt}}{k^s(1-\e^{-kt})},
       \qquad t>0,\quad0<a<1,\label{eq:F}\\
 \eps_m&=(-1)^{m+1},\qquad m\in\N,\quad m\ge1,\notag\\
 K_m(t,a)&=F_m(t,a)+\eps_mF_m(t,1-a),\label{eq:K}\\
 J_m(t,a)&=\left.\partial_s\left(F_s(t,a)
                    +\eps_mF_s(t,1-a)\right)\right|_{s=m}.
                    \label{eq:J}
\end{align}
The sign $\eps_m$ is held fixed during differentiation.
Let $P_m$ denote the explicit finite Laurent polynomial in
Theorem~\ref{thm:reflection}, and let $P_m^{[1]}$ denote the finite
power--logarithmic part in Proposition~\ref{prop:jet-tail}.
The main relations have the form
\begin{align}
 K_m&=P_m+N_m,\label{eq:overview-base}\\
 J_m&=P_m^{[1]}+t^m\Med\widetilde J_m+\Xi_m,\label{eq:overview-jet}\\
 t^m(\Sp-\Sm)\widetilde J_m&=\pi\ii N_m.\label{eq:overview-jump}
\end{align}
Here $\widetilde J_m$ is a Gevrey-1 formal series with an explicitly
meromorphic Borel transform. Both $N_m$ and $\Xi_m$ are convergent series
in $Q=\e^{-4\pi^2/t}$; the coefficients of $\Xi_m$ include $\log t$.
The linear median $\Med=(\Sp+\Sm)/2$ is defined only for the displayed
linear summation problem. No claim about nonlinear median operations
is needed.

\subsection{What is classical, and what is developed here}
The polylogarithm, Hurwitz-zeta special values, Bernoulli Fourier
expansions, and the zeta functional equations used below are classical
\cite{dlmf}. Dorigoni--Kleinschmidt analyze Lambert-series resurgence,
including the survival of nonperturbative terms when an odd-integer
asymptotic tail vanishes \cite[Section~2.2]{dk}. Thus neither the existence
of hidden exponential terms nor the general strategy of deforming a
parameter is being presented as a new phenomenon.

Fantini--Rella study $q$-Pochhammer resurgence and character-weighted
combinations \cite{fr}. A ProveIt report already gives an exact
Mellin--Borel completion for the order-one shifted $q$-factorial and a
finite Fourier--M\"obius criterion for rational weights \cite{repoq}.
Since $F_1(t,a)=-\log(\e^{-at};\e^{-t})_\infty$, that report directly
overlaps the $m=1$ background of this one.

The development here is the explicit first-order-jet completion for all
$m\ge1$ and all real $0<a<1$, the normalization-sensitive identity
\eqref{eq:overview-jump}, arbitrary-real-atomic reflection rigidity, its
sharp action-detection theorem, and the general-measure one-jet kernel.
The all-order collision formula supplies a separate bridge to
Hurwitz-zeta derivatives, harmonic numbers, and Stieltjes constants.
These are proved independently below. A limited literature comparison
cannot certify that every formula is absent from all earlier work.

\subsection{Repository boundary}
The detailed comparison with \cite{repoq} is pinned to
\begin{center}\small\pin.\end{center}
The Polylogarithms index, the Transseries index, and the incoming-directory
instructions were also inspected \cite{repoindex}. The binary incoming
archives were inventoried but not unpacked in this session. The
552-page combined Polylogarithms manuscript and the whole Transseries
collection were not audited. In particular, no claim is made to resolve
the separate S6/S8 polylogarithmic conjectures or to review every
incoming resonance manuscript. Section~\ref{sec:audit} records a narrow
cross-reference correction and distinguishes it from mathematical
corrections required by the new double-pole problem.

\section{Analytic setup and Mellin inversion}
\label{sec:mellin-setup}

We use $B_n(a)$ defined by
\begin{equation}
 \frac{z\e^{az}}{\e^z-1}=\sum_{n\ge0}B_n(a)\frac{z^n}{n!},
 \quad B_n=B_n(0),\quad B_1=-\tfrac12.
 \label{eq:Bernoulli}
\end{equation}
Write $H_n^{(r)}=\sum_{k=1}^nk^{-r}$, $H_n=H_n^{(1)}$,
and $H_0^{(r)}=0$. The Stieltjes convention is
\begin{equation}
 \zeta(1+u)=\frac1u+\sum_{j\ge0}\frac{(-1)^j\gamma_j}{j!}u^j,
 \qquad\gamma_0=\gamma.
 \label{eq:Stieltjes}
\end{equation}
A prime on $\zeta(w,a)$ always denotes differentiation in its first
argument, not in $a$.

\begin{lemma}[Analyticity and Mellin representation]\label{lem:Mellin}
For real $0<a<1$, $F_s(t,a)$ is entire in $s$ and holomorphic in
$\Re t>0$. For $t>0$ and $c>\max(1,1-\Re s)$,
\begin{equation}
 F_s(t,a)=\frac1{2\pi\ii}\int_{(c)}
  \Gamma(w)\zeta(w,a)\zeta(s+w)t^{-w}\dd w.
 \label{eq:MB}
\end{equation}
Every fixed order derivative is obtained by replacing $\zeta(s+w)$
with its corresponding derivative. These operations are uniform on
compact subsets of the stated domains, including real $a$-intervals
strictly inside $(0,1)$.
\end{lemma}
\begin{proof}
The second series in \eqref{eq:F} and each of its $s$-derivatives are
normally convergent on such compact sets: their terms are bounded by
$C\e^{-\delta k}k^M(1+\log k)^r$ for suitable $\delta>0$ and $M$.
Expanding $(1-\e^{-kt})^{-1}$ gives the first series and the absolutely
convergent double exponential sum. Insert
\[
 \e^{-x}=\frac1{2\pi\ii}\int_{(c)}\Gamma(w)x^{-w}\dd w
\]
into that sum. The two Dirichlet series converge absolutely when
$c>1$ and $c+\Re s>1$. Exponential vertical decay of $\Gamma$ justifies
Fubini and differentiation. This proves \eqref{eq:MB}.
\end{proof}

The standard identities needed later are
\begin{align}
 \zeta(-k,a)&=-\frac{B_{k+1}(a)}{k+1},\qquad k\ge0,
       \label{eq:zeta-B}\\
 B_n(1-a)&=(-1)^nB_n(a),\label{eq:B-reflection}\\
 \frac{B_n(a)}{n!}&=-\sum_{d\ne0}\frac{\e^{2\pi\ii da}}{(2\pi\ii d)^n},
       \qquad n\ge2,\label{eq:B-Fourier}\\
 \zeta(w,a)&=2\Gamma(1-w)(2\pi)^{w-1}
  \left(\sin\frac{\pi w}{2}\,C_{1-w}(a)
             +\cos\frac{\pi w}{2}\,S_{1-w}(a)\right),
       \quad\Re w<0,\label{eq:HFE}\\
 \zeta(z)&=2^z\pi^{z-1}\sin\frac{\pi z}{2}
                     \Gamma(1-z)\zeta(1-z),\label{eq:RFE}
\end{align}
where $C_v(a)=\sum_{d\ge1}\cos(2\pi da)d^{-v}$ and
$S_v(a)=\sum_{d\ge1}\sin(2\pi da)d^{-v}$ whenever these series converge
absolutely. The displayed applications of \eqref{eq:HFE} are on lines
where they do. Equations \eqref{eq:zeta-B}--\eqref{eq:RFE} are standard
\cite[Sections~24.8, 25.4, 25.11]{dlmf}.

Contour shifts to any fixed negative nonintegral line are justified by
vertical-strip bounds for the zeta functions and their derivatives,
together with the exponential factor from $\Gamma(w)$. They give
Poincar\'e expansions as $t\to0^+$, with a remainder estimated by the
new vertical integral. For complex $t$, the same argument works on
closed subsectors of $|\arg t|<\pi/2$. We do not obtain order derivatives
by differentiating a nonuniform, noninteger-order expansion: we first
differentiate the convergent Mellin integral.

\section{Collision jets and Stieltjes cancellation}
\label{sec:collision}

At $s=m\ge1$, the pole of $\zeta(s+w)$ meets the Gamma pole at
$w=1-m$. This collision controls the logarithmic part of the expansion.
The following coefficient formula handles every derivative order at
once.

\begin{theorem}[All-order collision residue]\label{thm:collision}
Set $n=m-1$, $L=\log t$, and
\begin{equation}
 B_{n,a}(u)=u\Gamma(-n+u)\zeta(-n+u,a).
 \label{eq:Bcollision}
\end{equation}
This function is holomorphic at $u=0$. The residue at $w=-n$ in the
Mellin representation of $\partial_s^rF_s(t,a)|_{s=m}$ is
\begin{equation}
 (-1)^rt^n\left\{
 r!\coef{u}{r+1}\bigl(\e^{-uL}B_{n,a}(u)\bigr)
                  +\gamma_r B_{n,a}(0)\right\},\qquad r\ge0.
 \label{eq:collision-general}
\end{equation}
In particular, this is a polynomial in $L$ of degree at most $r+1$.
\end{theorem}
\begin{proof}
The singular and constant terms of the order derivative are
\[
 \zeta^{(r)}(1+u)=\frac{(-1)^rr!}{u^{r+1}}
                            +(-1)^r\gamma_r+O(u).
\]
The remaining factors are $t^n\e^{-uL}B_{n,a}(u)/u$.
The singular term selects the coefficient of $u^{r+1}$, while the
constant term selects $B_{n,a}(0)$. Every $O(u)$ term has zero residue.
This proves \eqref{eq:collision-general}, including all signs and
factorials.
\end{proof}

For reflection, introduce
\begin{equation}
 Z_m(w,a)=\zeta(w,a)+\eps_m\zeta(w,1-a).
 \label{eq:Z}
\end{equation}
By \eqref{eq:zeta-B} and \eqref{eq:B-reflection},
$Z_m(1-m,a)=0$. Consequently the Gamma singularity in
\begin{equation}
 A_{m,a}(u)=\Gamma(1-m+u)Z_m(1-m+u,a)
 \label{eq:Acollision}
\end{equation}
is removable.

\begin{corollary}[Reflection removes one logarithmic degree]
\label{cor:collision-reflection}
For the fixed-parity order jet
$\partial_s^r(F_s(t,a)+\eps_mF_s(t,1-a))|_{s=m}$, the collision
contribution is exactly
\begin{equation}
 (-1)^rr!t^{m-1}\coef{u}{r}
                       \bigl(\e^{-uL}A_{m,a}(u)\bigr).
 \label{eq:collision-reflected}
\end{equation}
There is no explicit $\gamma_r B(0)$ term, and the logarithmic degree
is at most $r$. Its coefficient of $L^r$ is $t^{m-1}A_{m,a}(0)$.
\end{corollary}
\begin{proof}
For the reflected combination, the function in
\eqref{eq:Bcollision} is $uA_{m,a}(u)$, whose constant term is zero.
Substitute into Theorem~\ref{thm:collision}. The coefficient of $L^r$
comes from $(-uL)^r/r!$; its sign cancels the external $(-1)^r$.
\end{proof}

This cancellation refers to the explicit Stieltjes constants contributed
by the pole at $1$. It does not say that Hurwitz-zeta derivatives cannot
be reexpressed in bases involving other constants. It also does not
remove all logarithms: a first jet usually has a $t^{m-1}\log t$ term.

\subsection{Harmonic-number generation of the collision coefficients}
The finite Gamma factor has the convergent local expansion
\begin{equation}
 u\Gamma(-n+u)=\frac{(-1)^n}{n!}
 \exp\left((H_n-\gamma)u+
 \sum_{r\ge2}\frac{H_n^{(r)}+(-1)^r\zeta(r)}r u^r\right),
 \qquad |u|<1.
 \label{eq:Gamma-harmonic}
\end{equation}
Indeed, the left side equals
$(-1)^n\Gamma(1+u)\prod_{k=1}^n(1-u/k)^{-1}/n!$.
Taking its local logarithm proves \eqref{eq:Gamma-harmonic}. Together
with the Taylor coefficients of $\zeta(-n+u,a)$, this gives a finite,
algorithmic expression for every collision jet. Only finitely many
harmonic sums and Hurwitz-zeta derivatives are required at each order.

For example, with $Z_j=Z_m^{(j)}(1-m,a)$ and $h=H_{m-1}-\gamma$,
\begin{align}
 A_{m,a}(0)&=\frac{(-1)^{m-1}}{(m-1)!}Z_1,\notag\\
 A_{m,a}'(0)&=\frac{(-1)^{m-1}}{(m-1)!}
                               \left(\frac{Z_2}{2}+hZ_1\right),\notag\\
 \coef{u}{2}A_{m,a}(u)&=\frac{(-1)^{m-1}}{(m-1)!}
 \left(\frac{Z_3}{6}+\frac{hZ_2}{2}
       +\frac{h^2+H_{m-1}^{(2)}+\zeta(2)}2 Z_1\right).
 \label{eq:A-coefficients}
\end{align}
The second-order reflected collision is therefore
$t^{m-1}(A(0)L^2-2A'(0)L+A''(0))$.
This exact formula already reaches the Gamma/polygamma and Stieltjes
part of the Polylogarithms program without conjectural constant
reduction.

\section{The ordinary integer-order Borel kernel}
\label{sec:base-borel}

Before imposing reflection, the integer-order expansion of a single
shift is
\begin{equation}
 F_m(t,a)\sim P_m^F(t,a)+t^{m-1}\widetilde R_{m,a}(t),
 \qquad\widetilde R_{m,a}(t)=\sum_{j\ge1}c_{m,j}(a)t^{2j},
 \label{eq:F-expansion}
\end{equation}
where
\begin{align}
 P_m^F(t,a)&=\frac{\zeta(m+1)}t+
 \sum_{\substack{0\le k\le m\\k\ne m-1}}
       \frac{(-t)^k}{k!}\zeta(-k,a)\zeta(m-k)\notag\\
 &\quad+\frac{(-t)^{m-1}}{(m-1)!}
 \left(\zeta'(1-m,a)
       +(H_{m-1}-\log t)\zeta(1-m,a)\right),\label{eq:PF}\\
 c_{m,j}(a)&=\frac{(-1)^{m-1}B_{m+2j}(a)B_{2j}}
                         {(m+2j)!\,2j}.
                         \label{eq:c-coeff}
\end{align}
These formulas follow by shifting \eqref{eq:MB}. The pole at $w=1$
gives the $t^{-1}$ term, the collision is Theorem~\ref{thm:collision}
with $r=0$, and the trivial zeta zeros leave only the displayed tail.

\subsection{Normalization and continuation}
For a formal power series with zero constant term we use
\begin{equation}
 \B(t^r)=\frac{\xi^{r-1}}{(r-1)!},\qquad
 \mathcal L_\theta b(t)=\int_0^{\e^{\ii\theta}\infty}
                                      \e^{-\xi/t}b(\xi)\dd\xi.
 \label{eq:Borel-normalization}
\end{equation}
There is no factor $1/t$ in the Laplace transform. Let
\begin{align}
 A_N&=4\pi^2N,\label{eq:actions}\\
 D_m(N,a)&=\sum_{d\mid N}
       \frac{\e^{2\pi\ii da}+(-1)^m\e^{-2\pi\ii da}}{d^m}.
       \label{eq:D}
\end{align}

\begin{theorem}[Meromorphic simple-pole kernel]\label{thm:base-borel}
The Borel transform of $\widetilde R_{m,a}$ is
\begin{equation}
 \B\widetilde R_{m,a}(\xi)=
 \frac{2(-1)^{m-1}}{(2\pi\ii)^m}
 \sum_{N\ge1}D_m(N,a)\frac{\xi}{A_N^2-\xi^2}.
 \label{eq:base-kernel}
\end{equation}
The series is normally convergent off the lattice $\{\pm A_N:N\ge1\}$.
At either $\xi=A_N$ or $\xi=-A_N$, its residue is
\begin{equation}
 \frac{(-1)^m}{(2\pi\ii)^m}D_m(N,a).
 \label{eq:base-residue}
\end{equation}
A displayed pole is removable exactly when this coefficient vanishes.
The formal tail is Gevrey-1 and is Borel--Laplace summable in every
nonreal direction that is compatible with the Laplace exponential.
\end{theorem}
\begin{proof}
Insert \eqref{eq:B-Fourier} and
$B_{2j}=(-1)^{j+1}2(2j)!\zeta(2j)/(2\pi)^{2j}$ into
\eqref{eq:c-coeff}. Expand $\zeta(2j)$ as a Dirichlet series, pair the
positive and negative Fourier indices, and sum
\[
 \sum_{j\ge1}\frac{\xi^{2j-1}}{(4\pi^2dn)^{2j}}
                  =\frac{\xi}{(4\pi^2dn)^2-\xi^2}.
\]
This gives \eqref{eq:base-kernel} near zero, with $N=dn$.
For every $N$, $|D_m(N,a)|\le2\sum_{d\mid N}d^{-m}
\le2(1+\log N)$. On a compact set, the large-$N$ summands are
$O((1+\log N)N^{-2})$, proving normal convergence and meromorphic
continuation. Each rational summand has residue $-1/2$ at both of its
poles, which gives \eqref{eq:base-residue}.

On a ray separated from the real axis,
$|A_N^2-\xi^2|\ge c_\theta(A_N^2+|\xi|^2)$.
Split the sum at $N\asymp1+|\xi|$ to obtain at most logarithmic growth
at infinity, and $O(\xi)$ at zero. Thus the corresponding Laplace
integrals converge, and their termwise evaluation is justified.
Analyticity of the Borel germ in $|\xi|<4\pi^2$ gives the Gevrey-1
coefficient estimate by Cauchy's inequality. Expanding the integrand
near zero and estimating the remainder of the Laplace integral gives
the stated formal expansion of its sum.
\end{proof}

Reflection removes this entire kernel, since its Bernoulli coefficient
has degree $m+2j$ and
$B_{m+2j}(a)+\eps_m B_{m+2j}(1-a)=0$.
It is essential not to conclude from this that the reflected analytic
function equals its finite polynomial.

\section{Exact reflection and a convergent dual polylogarithm}
\label{sec:reflection}

Define
\begin{align}
 h_m(d,a)&=\e^{2\pi\ii da}+\eps_m\e^{-2\pi\ii da},\notag\\
 E_m(N,a)&=\sum_{d\mid N}\frac{h_m(d,a)}{d^m},\label{eq:E}\\
 \Phi_m(Q,a)&=\sum_{N\ge1}E_m(N,a)Q^N\notag\\
 &=\sum_{\ell\ge1}\left(\Li_m(\e^{2\pi\ii a}Q^\ell)
              +\eps_m\Li_m(\e^{-2\pi\ii a}Q^\ell)\right),
       \quad |Q|<1.\label{eq:Phi}
\end{align}
These are normally convergent power series. The factor $(-1)^m$ in
$D_m$ has changed to $\eps_m=(-1)^{m+1}$ in $E_m$: the two divisor
sequences encode complementary reflection components.

\begin{theorem}[Exact reflected duality]\label{thm:reflection}
For every integer $m\ge1$, real $0<a<1$, and $t>0$,
\begin{equation}
 K_m(t,a)=P_m(t,a)+\left(\frac{\ii t}{2\pi}\right)^{m-1}
                         \Phi_m(\e^{-4\pi^2/t},a),
 \label{eq:reflection-exact}
\end{equation}
where the finite Laurent polynomial is
\begin{align}
 P_m(t,a)&=\frac{(1+\eps_m)\zeta(m+1)}t
 +\sum_{\substack{0\le k\le m\\k\ne m-1}}
       \frac{(-t)^k}{k!}\zeta(m-k)Z_m(-k,a)\notag\\
 &\quad+\frac{(-t)^{m-1}}{(m-1)!}Z_m'(1-m,a).
 \label{eq:P-reflection}
\end{align}
In particular, $N_m=K_m-P_m$ is given exactly by the convergent dual
series, not merely by its first exponential term.
\end{theorem}
\begin{proof}
In the Mellin integral for $K_m$, the zeta factor $\zeta(w,a)$ is replaced
by $Z_m(w,a)$. Shift to $\Re w=-m-\delta$, where $0<\delta<1$.
The crossed residues are precisely \eqref{eq:P-reflection}: the
collision loses its logarithm because $Z_m(1-m,a)=0$.

Put $v=1-m-w$ and
\[
 H_m(v,a)=\sum_{d\ge1}\frac{h_m(d,a)}{d^{m+v}}.
\]
Combining \eqref{eq:HFE}, \eqref{eq:RFE}, and
$\Gamma(w)\Gamma(1-w)=\pi/\sin(\pi w)$ gives
\begin{equation}
 \Gamma(w)\zeta(m+w)Z_m(w,a)
 =\left(\frac{\ii}{2\pi}\right)^{m-1}(4\pi^2)^{-v}
              \Gamma(v)\zeta(v)H_m(v,a).
 \label{eq:FE-product}
\end{equation}
For clarity, when $m$ is odd, $Z_m$ keeps the cosine Fourier series;
when $m$ is even, it keeps the sine Fourier series, with
$H_m=2\ii S_{m+v}$. In either case the elementary trigonometric factors
in \eqref{eq:FE-product} reduce to the displayed constant. The identity
holds across any removable zeros by analytic continuation.

The reflected remainder is consequently
\begin{equation}
 \left(\frac{\ii t}{2\pi}\right)^{m-1}
 \frac1{2\pi\ii}\int_{(1+\delta)}
  \Gamma(v)\zeta(v)H_m(v,a)X^{-v}\dd v,
 \qquad X=\frac{4\pi^2}t.
 \label{eq:reflection-MB}
\end{equation}
Both Dirichlet series converge absolutely on this line. Termwise
Mellin inversion evaluates \eqref{eq:reflection-MB} as
\[
 \left(\frac{\ii t}{2\pi}\right)^{m-1}
 \sum_{d,n\ge1}\frac{h_m(d,a)}{d^m}\e^{-Xdn},
\]
which is \eqref{eq:reflection-exact}. All sums in this last expression
converge absolutely for every $t>0$.
\end{proof}

\begin{remark}[An exact transseries, not only an asymptotic statement]
The action set is $4\pi^2\N$. It is locally finite, and the exponential
block is an ordinary convergent power series in $Q$. Its coefficients
are explicit divisor sums. The identity therefore supplies an analytic
realization of the full exponentially small block, while the ordinary
power expansion contains none of that block's coefficients.
\end{remark}

\section{The order derivative and its double-pole kernel}
\label{sec:jet-borel}

\begin{proposition}[Finite part and activated coefficients]
\label{prop:jet-tail}
The fixed-parity first jet \eqref{eq:J} has the expansion
\begin{equation}
 J_m(t,a)\sim P_m^{[1]}(t,a)+t^m\widetilde J_m(t,a),
 \qquad\widetilde J_m(t,a)=\sum_{j\ge1}d_{m,j}(a)t^{2j},
 \label{eq:jet-expansion}
\end{equation}
where
\begin{align}
 P_m^{[1]}(t,a)&=\frac{(1+\eps_m)\zeta'(m+1)}t+
 \sum_{\substack{0\le k\le m\\k\ne m-1}}
       \frac{(-t)^k}{k!}\zeta'(m-k)Z_m(-k,a)\notag\\
 &\quad+\frac{(-t)^{m-1}}{(m-1)!}
 \left((\log t-H_{m-1}+\gamma)Z_m'(1-m,a)
                          -\frac12Z_m''(1-m,a)\right),
       \label{eq:P-jet}\\
 d_{m,j}(a)&=\frac{(-1)^{m+j+1}(2j)!\,B_{m+2j+1}(a)\zeta(2j+1)}
                        {(m+2j+1)!(2\pi)^{2j}}.
       \label{eq:d-coeff}
\end{align}
\end{proposition}
\begin{proof}
Differentiate the Mellin integral first and then shift its contour.
The finite residues away from $w=1-m$ give the first line of
\eqref{eq:P-jet}. Corollary~\ref{cor:collision-reflection} with $r=1$
gives the second line.

For $k>m$, the coefficient at $t^k$ is
$(-1)^kZ_m(-k,a)\zeta'(m-k)/k!$.
When $k=m+2j-1$, the Bernoulli reflection makes $Z_m(-k,a)=0$.
When $k=m+2j$, the zero of $\zeta(-2j)$ is replaced by its nonzero
derivative. Differentiating \eqref{eq:RFE} at $-2j$ gives
\begin{equation}
 \zeta'(-2j)=\frac{(-1)^j(2j)!}{2(2\pi)^{2j}}\zeta(2j+1),
 \qquad j\ge1.
 \label{eq:trivial-zero-derivative}
\end{equation}
Also $Z_m(-m-2j,a)=-2B_{m+2j+1}(a)/(m+2j+1)$.
Substitution proves \eqref{eq:d-coeff}.
\end{proof}

The entire activated tail is therefore driven by derivatives at the
trivial zeros of the Riemann zeta function. This explains both its odd
zeta values and its one-step change of power parity.

\begin{theorem}[Explicit double-pole Borel transform]
\label{thm:jet-borel}
Let $C_m=(-1)^m/(2\pi\ii)^{m+1}$. Then
\begin{equation}
 \B\widetilde J_m(\xi,a)=C_m\sum_{N\ge1}\frac{E_m(N,a)}N
                  \frac{2A_N^2\xi}{(A_N^2-\xi^2)^2}.
 \label{eq:jet-kernel}
\end{equation}
This is normally convergent off the action lattice and is of at most
polynomial growth on rays separated from the real axis. At $\xi=A_N$,
the coefficient of $(\xi-A_N)^{-2}$ is
\begin{equation}
 2\pi^2 C_mE_m(N,a),
 \label{eq:double-coeff}
\end{equation}
and the coefficient of $(\xi-A_N)^{-1}$ is zero. The double-pole
coefficient at $-A_N$ is the negative of \eqref{eq:double-coeff}.
The series $\widetilde J_m$ is Gevrey-1 and summable off the real Borel
rays.
\end{theorem}
\begin{proof}
Apply the Bernoulli Fourier expansion to \eqref{eq:d-coeff} and expand
$\zeta(2j+1)$. After division by $(2j-1)!$, the extra factor is $2j$.
For $A\ne0$,
\begin{equation}
 \sum_{j\ge1}2j\frac{\xi^{2j-1}}{A^{2j}}
                  =\frac{2A^2\xi}{(A^2-\xi^2)^2}.
 \label{eq:double-geometric}
\end{equation}
The leftover divisor weight is $d^{-m}/N$, giving
\eqref{eq:jet-kernel} with the stated constant.
The large-$N$ bound on compact sets is now
$O((1+\log N)N^{-3})$. On an allowed ray, use
$|A_N^2-\xi^2|\ge c(A_N^2+|\xi|^2)$ and split at
$N\asymp1+|\xi|$; this proves polynomial growth and Laplace convergence.

Finally,
\begin{equation}
 \frac{2A^2\xi}{(A^2-\xi^2)^2}
          =\frac{\dd}{\dd\xi}\frac{\xi^2}{A^2-\xi^2}.
 \label{eq:kernel-derivative}
\end{equation}
A derivative of a meromorphic function has zero residue. Its double
coefficients at $A$ and $-A$ are $A/2$ and $-A/2$, respectively.
Multiplication by $C_mE_m(N,a)/N$ proves \eqref{eq:double-coeff}.
The same Cauchy and Laplace arguments as in
Theorem~\ref{thm:base-borel} give the Gevrey and summability assertions.
\end{proof}

A zero residue of the \emph{Borel kernel} does not mean a zero Stokes
jump. The residue must be taken after multiplication by the Laplace
exponential; that operation turns each double pole into a nonzero
simple-pole coefficient. The next section computes it with its sign.

\section{Exact order-jet completion and recovery of the hidden block}
\label{sec:completion}

For $t>0$, define $\Sp$ and $\Sm$ by Laplace rays immediately above
and below the positive Borel axis, respectively. More precisely, one
can use fixed rays of angles $\theta$ and $-\theta$, where
$0<\theta<\pi/2$, and deform within the indicated half-planes. Their
values are independent of the particular such angle. The difference
of the upper and lower contours is clockwise around positive poles.

\subsection{The finite-part kernel needed for a double pole}

\begin{lemma}[A simple-pole representation of the linear median]
\label{lem:I}
For $\lambda>0$, let
\begin{equation}
 I(\lambda)=\PV\int_0^\infty \e^{-\lambda y}\frac{y^2}{1-y^2}\dd y
 =-\frac1\lambda+
 \frac12\left(\e^{-\lambda}\operatorname{Ei}(\lambda)
                    +\e^\lambda E_1(\lambda)\right).
 \label{eq:I}
\end{equation}
It is $O(\lambda^{-1}+1)$ at zero and $O(\lambda^{-3})$ at infinity.
For $1<\Re v<3$,
\begin{equation}
 \int_0^\infty\lambda^{v-1}I(\lambda)\dd\lambda
                      =\frac\pi2\Gamma(v)\tan\frac{\pi v}{2}.
 \label{eq:I-Mellin}
\end{equation}
For $A>0$, $t>0$, and $\lambda=A/t$, the upper and lower Laplace
integrals of the double-pole kernel satisfy
\begin{equation}
 \int_0^{\e^{\pm\ii\theta}\infty}
 \e^{-\xi/t}\frac{2A^2\xi}{(A^2-\xi^2)^2}\dd\xi
 =\lambda I(\lambda)\ \pm\ \frac{\pi\ii\lambda}{2}\e^{-\lambda}.
 \label{eq:double-Laplace}
\end{equation}
\end{lemma}
\begin{proof}
Use $y^2/(1-y^2)=-1+1/(1-y^2)$ and partial fractions. The elementary
exponential-integral evaluations give \eqref{eq:I}. They give the
endpoint estimates as well: in particular
$I(\lambda)\sim\sum_{j\ge1}(2j)!\lambda^{-2j-1}$ at infinity.
This last expansion is asymptotic, not a convergent series.

For the Mellin transform, first omit $(1-\delta,1+\delta)$ from the
$y$-integral. Away from the pole, the absolute integral after integrating
in $\lambda$ converges at both endpoints precisely when
$1<\Re v<3$. Subtracting the value of the pole numerator at $y=1$
shows that the difference between the truncated and principal-value
kernels is bounded by
$C\delta(1+\lambda)\e^{-\lambda/2}$ for $0<\delta<1/2$.
This bound is integrable against $\lambda^{\Re v-1}\dd\lambda$.
Consequently the limit commutes with the Mellin integral, giving
\begin{align*}
 \Gamma(v)\PV\int_0^\infty\frac{y^{2-v}}{1-y^2}\dd y
 &=\frac{\Gamma(v)}2\PV\int_0^\infty\frac{u^{(1-v)/2}}{1-u}\dd u\\
 &=\frac\pi2\Gamma(v)\cot\frac{\pi(3-v)}2
  =\frac\pi2\Gamma(v)\tan\frac{\pi v}2.
\end{align*}
The principal-value cotangent identity follows by splitting at $1$,
inverting the integration variable in the second part, and using the
partial fraction expansion of the cotangent. Mellin inversion is
valid on any vertical line $1<c<3$; the Gamma factor supplies
exponential decay there.

For \eqref{eq:double-Laplace}, use \eqref{eq:kernel-derivative} and
integrate by parts on a nonreal ray. The endpoint contributions vanish.
After $\xi=Ay$, the resulting integral is $\lambda$ times the lateral
integral defining $I$. The residue of $y^2/(1-y^2)$ at $y=1$ is $-1/2$.
The upper contour contributes $+\pi\ii\e^{-\lambda}/2$ relative to
the symmetric principal value, and the lower contour contributes its
negative. This proves the formula and fixes its orientation.
\end{proof}

The double-pole Laplace integral itself does \emph{not} possess an
ordinary Cauchy principal value on the real line. Its linear median is
a Hadamard finite part, equivalently the simple-pole expression in
\eqref{eq:double-Laplace}. This distinction is necessary even though
the double-pole kernel has zero residue before the exponential is
inserted.

\subsection{The exact logarithmic dual correction}
Define
\begin{align}
 \Xi_m(t,a)&=\left(\frac{\ii t}{2\pi}\right)^{m-1}
 \sum_{d\ge1}\frac{h_m(d,a)}{d^m}
             \frac{\log(t/(2\pi d))}{\e^{4\pi^2d/t}-1}.
 \label{eq:Xi}
\end{align}
For every $t>0$ this series converges absolutely. It is an exact
polylogarithmic order-derivative block: if
\begin{equation}
 \Phi_{s\mid m}(Q,a)=\sum_{\ell\ge1}
 \left(\Li_s(\e^{2\pi\ii a}Q^\ell)
              +\eps_m\Li_s(\e^{-2\pi\ii a}Q^\ell)\right),
 \label{eq:fixed-Phi}
\end{equation}
then
\begin{equation}
 \Xi_m(t,a)=\left(\frac{\ii t}{2\pi}\right)^{m-1}
 \left(\log\frac{t}{2\pi}\,\Phi_m(Q,a)
          +\left.\partial_s\Phi_{s\mid m}(Q,a)\right|_{s=m}\right).
 \label{eq:Xi-order}
\end{equation}
The parity sign, again, is fixed.

\begin{theorem}[Exact first-jet completion and activation]
\label{thm:activation}
Under the hypotheses of Theorem~\ref{thm:reflection},
\begin{align}
 J_m(t,a)&=P_m^{[1]}(t,a)+t^m\Med\widetilde J_m(t,a)+\Xi_m(t,a),
                    \label{eq:exact-J-median}\\
 t^m(\Sp-\Sm)\widetilde J_m(t,a)
              &=\pi\ii\bigl(K_m(t,a)-P_m(t,a)\bigr).
                    \label{eq:activation}\\
 J_m(t,a)&=P_m^{[1]}(t,a)+t^m\mathcal S_\pm\widetilde J_m(t,a)
             +\Xi_m(t,a)\mp\frac{\pi\ii}{2}N_m(t,a).
                    \label{eq:exact-J-lateral}
\end{align}
The linear median also has the absolutely convergent kernel formula
\begin{equation}
 t^m\Med\widetilde J_m(t,a)
       =\left(\frac{\ii t}{2\pi}\right)^{m-1}
                         \sum_{N\ge1}E_m(N,a)I(A_N/t).
 \label{eq:median-I}
\end{equation}
\end{theorem}
\begin{proof}
\emph{Identification of the analytic function.}
Differentiate \eqref{eq:MB} in $s$ before specializing to $s=m$, and
then form the reflected combination. Shift to
$\Re w=-m-\delta$ with $0<\delta<1$. Proposition~\ref{prop:jet-tail}
computes the crossed residues as $P_m^{[1]}$.

In the remainder set $v=1-m-w$ as in
Theorem~\ref{thm:reflection}. Logarithmic differentiation of the
Riemann functional equation, wherever none of its factors is zero,
gives the multiplier
\[
 \log(2\pi)-\psi(v)-\frac{\zeta'(v)}{\zeta(v)}
                                +\frac\pi2\tan\frac{\pi v}{2}.
\]
To avoid any artificial division by a zeta zero, put
$G(v)=\Gamma(v)\zeta(v)$ and use the holomorphic product expression
\begin{equation}
 \log(2\pi)G(v)-G'(v)+\frac\pi2\tan\frac{\pi v}{2}G(v).
 \label{eq:FE-jet-product}
\end{equation}
The identity extends across removable zeros. The exact remainder is
\begin{equation}
 \left(\frac{\ii t}{2\pi}\right)^{m-1}
 \frac1{2\pi\ii}\int_{(1+\delta)}
 \left(\log(2\pi)G-G'+\frac\pi2\tan\frac{\pi v}2G\right)
                         H_m(v,a)X^{-v}\dd v.
 \label{eq:jet-remainder-MB}
\end{equation}
There are no tangent poles on this line.

Integrate the $-G'$ term by parts along the vertical contour. The
horizontal endpoint terms vanish by Gamma decay. The first two terms
of the integrand become
\[
 G(v)\left(\log\frac{t}{2\pi}H_m(v,a)+H_m'(v,a)\right)X^{-v},
\]
since $\log(2\pi)-\log X=\log(t/(2\pi))$.
The series for $H_m$ and $H_m'$ converge absolutely. Mellin inversion
of their Dirichlet expansions therefore gives exactly \eqref{eq:Xi}.

For the last term in \eqref{eq:jet-remainder-MB}, expand
$\zeta(v)H_m(v,a)=\sum_{N\ge1}E_m(N,a)N^{-v}$.
Lemma~\ref{lem:I} identifies the resulting inverse Mellin integral with
$I(XN)$. Thus the remainder not accounted for by $\Xi_m$ is the
right side of \eqref{eq:median-I}.
This series converges absolutely because
$E_m(N,a)=O(1+\log N)$ and $I(XN)=O(N^{-3})$ for fixed $X>0$.

\emph{Identification with the Borel median.}
Integrate \eqref{eq:jet-kernel} termwise on upper and lower rays,
which is justified by the absolute majorants in
Theorem~\ref{thm:jet-borel}. Lemma~\ref{lem:I} gives their average.
The constants simplify using
\begin{equation}
 4\pi^2 C_m=\left(\frac{\ii}{2\pi}\right)^{m-1}.
 \label{eq:constant-check}
\end{equation}
Indeed, $A_N/(tN)=4\pi^2/t$ in every summand. This proves
\eqref{eq:median-I}, and hence \eqref{eq:exact-J-median}.

\emph{The discontinuity.}
By \eqref{eq:double-Laplace}, the difference of the two lateral
integrals of one rational summand is
$\pi\ii(A_N/t)\e^{-A_N/t}$. Multiplying by
$t^mC_mE_m(N,a)/N$ and summing gives
\[
 4\pi^3\ii C_m t^{m-1}\sum_{N\ge1}E_m(N,a)Q^N
 =\pi\ii\left(\frac{\ii t}{2\pi}\right)^{m-1}\Phi_m(Q,a).
\]
Theorem~\ref{thm:reflection} identifies the last factor with $N_m$.
This proves \eqref{eq:activation}. Combining it with the median identity
gives \eqref{eq:exact-J-lateral}.
\end{proof}

\begin{remark}[A direct residue check]
At $\xi=A_N$, write the Borel principal part as
$c_N(\xi-A_N)^{-2}$, with $c_N=2\pi^2C_mE_m(N,a)$.
The residue of the \emph{Laplace integrand} is
$-c_Nt^{-1}\e^{-A_N/t}$. The upper-minus-lower contour multiplies this
by $-2\pi\ii$. This reproduces the positive factor in
\eqref{eq:activation} and independently explains why simply taking
residues of the unweighted Borel kernel would miss the jump.
\end{remark}

\subsection{Complex sectors and transseries interpretation}
All exact identities proved for $t>0$ continue to $\Re t>0$ with the
principal $\Log t$ and the corresponding analytic continuations of the
lateral sums. To see this without manipulating formal symbols,
$I(A_N/t)$ is holomorphic there by \eqref{eq:I}; its large-$N$
estimates are uniform on compact subsets. The sums in
\eqref{eq:median-I}, \eqref{eq:Phi}, and \eqref{eq:Xi} are consequently
normally convergent. Both sides of the identities are holomorphic and
agree on the positive real axis. Alternatively, fixed nonreal Laplace
rays give overlapping sectorial definitions of the same continuations.

On a closed subsector $|\arg t|\le\pi/2-\eta$,
$\Re(4\pi^2/t)\ge4\pi^2\sin\eta/|t|$. The factors $Q^N$ are then
uniformly exponentially small. The exact structure of the first jet is
one Borel-summed power block plus convergent sectors
$t^{m-1}(\alpha_N\Log t+\beta_N)Q^N$.
No accumulating action set or formal-support completion is needed for
this particular transseries. The boundary directions $\arg t=\pm\pi/2$
are not covered by this small-exponential assertion.

\section{Finite reflection rigidity and action detection}
\label{sec:atomic}

Let
\begin{equation}
 \mu=\sum_{j=1}^r c_j\delta_{a_j},\qquad
 0<a_j<1,
 \label{eq:atomic}
\end{equation}
where the $a_j$ are distinct and the displayed $c_j$ are nonzero.
Let $R(a)=1-a$, write $R_*\mu$ for the reflected measure, and define
\begin{equation}
 F_{s,\mu}(t)=\int F_s(t,a)\dd\mu(a),\qquad
 \muh(n)=\int\e^{2\pi\ii na}\dd\mu(a).
 \label{eq:weighted}
\end{equation}
The atoms need not be rational. By a \emph{terminating tail} we mean
that the integer-order power--logarithmic expansion has only finitely
many nonzero coefficients. The fixed finite logarithmic collision part
is not included in the word ``tail.''

\begin{theorem}[Atomic reflection rigidity with a finite certificate]
\label{thm:atomic}
For $m\ge1$, the following conditions are equivalent:
\begin{enumerate}[label=(\roman*),leftmargin=12mm]
\item the asymptotic expansion of $F_{m,\mu}$ terminates;
\item all its tail coefficients $t^{m+2j-1}$, $j\ge1$, vanish;
\item $R_*\mu=\eps_m\mu$;
\item all positive-pole coefficients in its normalized Borel tail vanish.
\end{enumerate}
Let $M$ be the number of distinct points in
$\operatorname{supp}(\mu)\cup R(\operatorname{supp}(\mu))$.
For condition (iv), it suffices to check the first $M$ action locations;
in particular, at most $2r$ checks suffice.
\end{theorem}
\begin{proof}
The weighted Borel kernel in Theorem~\ref{thm:base-borel} has divisor
coefficients
\begin{equation}
 D_{m,\mu}(N)=\sum_{d\mid N}\frac{h_\mu(d)}{d^m},
 \quad h_\mu(d)=\muh(d)+(-1)^m\muh(-d).
 \label{eq:D-mu}
\end{equation}
M\"obius inversion is the exact identity
\begin{equation}
 h_\mu(n)=n^m\sum_{d\mid n}\mu_{\mathrm{Mob}}(n/d)D_{m,\mu}(d).
 \label{eq:Mobius}
\end{equation}
Thus the vanishing of $D_{m,\mu}(1),\ldots,D_{m,\mu}(M)$ implies
$h_\mu(1)=\cdots=h_\mu(M)=0$.

The latter are the first $M$ Fourier coefficients of
$\nu=\mu+(-1)^mR_*\mu$. Write this measure on its at most $M$
distinct nodes, with $z_j=\e^{2\pi\ii a_j}$. The matrix
$(z_j^n)_{1\le n,j\le M}$ is invertible: its determinant is a nonzero
product of the $z_j$ and the Vandermonde factors $z_k-z_j$.
Padding with zero coefficients if necessary gives $\nu=0$, which is
condition (iii).

Conversely, (iii) and Bernoulli reflection give (ii), and (ii) gives
(i). If (i) holds, the Borel germ is a polynomial. Its meromorphic
continuation \eqref{eq:base-kernel} must equal that polynomial wherever
both are defined; hence every residue vanishes and (iv) holds.
This proves all implications and the finite certificate.
\end{proof}

The finite certificate refers to action-pole data, not to an arbitrary
list of the first $M$ Taylor coefficients. Those are different finite
linear systems. At rational shifts, finite Fourier transforms give a
natural alternative presentation; for arbitrary real shifts the
Vandermonde proof avoids any rationality hypothesis.

\begin{theorem}[Hidden actions are visible in one order derivative]
\label{thm:action-detection}
Suppose the nonzero measure \eqref{eq:atomic} satisfies
$R_*\mu=\eps_m\mu$. Then its exact hidden remainder is
\begin{equation}
 N_{m,\mu}(t)=F_{m,\mu}(t)-P_{m,\mu}(t)
 =\left(\frac{\ii t}{2\pi}\right)^{m-1}
       \sum_{N\ge1}b_{m,\mu}(N)Q^N,
 \quad b_{m,\mu}(N)=\sum_{d\mid N}\frac{\muh(d)}{d^m}.
 \label{eq:N-mu}
\end{equation}
Here $P_{m,\mu}=\tfrac12\int P_m(t,a)\dd\mu(a)$.
The first nonzero coefficient in \eqref{eq:N-mu} occurs at an index
$N_0\le r$. The first order derivative has its first nonzero positive
double Borel pole at exactly $A_{N_0}$, and its jump is
$\pi\ii N_{m,\mu}$.
\end{theorem}
\begin{proof}
Since $R_*\mu=\eps_m\mu$, integration of a reflected pair gives
$\int K_m(t,a)\dd\mu(a)=2F_{m,\mu}(t)$. Also
$\muh(-d)=\eps_m\muh(d)$, so that
$\tfrac12\int E_m(N,a)\dd\mu(a)=b_{m,\mu}(N)$.
Integrating Theorem~\ref{thm:reflection} proves \eqref{eq:N-mu}.
The same reasoning applies to the fixed-parity order derivative;
Theorems~\ref{thm:jet-borel} and \ref{thm:activation} identify its
double-pole coefficient and its discontinuity.

If $\muh(1),\ldots,\muh(r)$ all vanished, the $r\times r$
Vandermonde system for the original atomic measure would force
$c_1=\cdots=c_r=0$. Choose the first nonzero Fourier coefficient,
with index $n_0\le r$. Every proper-divisor term in
$b_{m,\mu}(n_0)$ vanishes, leaving
$b_{m,\mu}(n_0)=\muh(n_0)/n_0^m\ne0$. Earlier $b$'s are zero.
This proves $N_0=n_0\le r$ and the matching double-pole assertion.
\end{proof}

On prescribed nodes, the hidden coefficients recover the weights by
an explicit two-stage linear procedure. First use
\begin{equation}
 \muh(n)=n^m\sum_{d\mid n}\mu_{\mathrm{Mob}}(n/d)b_{m,\mu}(d).
 \label{eq:recover-muh}
\end{equation}
Then solve the Vandermonde system on the known $\e^{2\pi\ii a_j}$.
This is an exact identification statement; it does not assert numerical
stability for nearly colliding nodes.

\subsection{Sharpness by midpoint grids}
For $r\ge1$, set
\begin{equation}
 \mu_r=\sum_{j=0}^{r-1}\delta_{(j+1/2)/r}.
 \label{eq:midpoint}
\end{equation}
This measure is reflection-symmetric. Its Fourier coefficients are
\begin{equation}
 \widehat\mu_r(n)=
 \begin{cases}
 r(-1)^k,&n=rk,\\
 0,&r\nmid n.
 \end{cases}
 \label{eq:midpoint-Fourier}
\end{equation}
For every odd $m$, Theorem~\ref{thm:action-detection} therefore has
$N_0=r$. In particular,
\begin{equation}
 N_{m,\mu_r}(t)=
 -r^{1-m}\left(\frac{\ii t}{2\pi}\right)^{m-1}
 \e^{-4\pi^2r/t}+O\left(t^{m-1}\e^{-8\pi^2r/t}\right)
 \quad(t\to0^+).
 \label{eq:midpoint-first}
\end{equation}
The first $r-1$ activated poles vanish, so the bound $r$ cannot be
improved uniformly. This arithmetic cancellation also has the exact
lattice explanation
\begin{equation}
 \sum_{j=0}^{r-1}F_s\left(t,\frac{j+1/2}{r}\right)
                         =F_s(t/r,1/2),
 \label{eq:grid-distribution}
\end{equation}
obtained by partitioning the index $r\ell+j$. Equation
\eqref{eq:grid-distribution} holds for every complex $s$, not only at
integer order, and can therefore be differentiated in $s$.

\section{The one-jet kernel for general measures}
\label{sec:general-measures}

Atomicity excludes a genuine continuous blind component. Let $\lambda$
denote Lebesgue measure of mass one on $(0,1)$. Direct integration gives
\begin{equation}
 \int_0^1 F_s(t,a)\dd a=\frac{\zeta(s+1)}t,
 \qquad\Re s>0.
 \label{eq:Haar}
\end{equation}
Thus a nonzero continuous measure can have both its base tail and every
order-derivative tail equal to zero. The next theorem shows that this
is the \emph{only} simultaneous blind component.

For a finite complex Borel measure $\mu$ on $(0,1)$, its Bernoulli
moments $M_n(\mu)=\int B_n(a)\dd\mu(a)$ always exist. Formal high-tail
coefficients are defined by integrating the polynomial coefficients of
the single-shift expansion. For an analytic interpretation of both
$F_{m,\mu}$ and its first jet, it suffices to assume
\begin{equation}
 \int_{(0,1)}\left(1+|\log a|^2+|\log(1-a)|^2\right)\dd|\mu|(a)<\infty.
 \label{eq:measure-integrability}
\end{equation}
Indeed, for $m=1$ the original series and its first jet grow at most
like $1+|\log a|$ and $1+|\log a|^2$ near the singular endpoint;
these bounds follow by splitting their $k$-sums at $k\asymp1/a$.
For $m\ge2$ the corresponding endpoint bounds are no worse.
The coefficient theorem below does not require interchanging a measure integral with the initial Mellin contour; such an interchange would require stronger endpoint moment hypotheses.

\begin{theorem}[Complete one-jet kernel]\label{thm:one-jet}
Fix $m\ge1$. On finite complex Borel measures on $(0,1)$, the joint
kernel of the two maps
\begin{equation}
 \mu\longmapsto\text{high tail of }F_{m,\mu},\qquad
 \mu\longmapsto\text{high tail of }
                          \left.\partial_sF_{s,\mu}\right|_{s=m}
 \label{eq:jet-map}
\end{equation}
is exactly $\C\lambda$. The same conclusion holds when ``zero high
tail'' is replaced by ``eventually terminating high tail.''
Consequently the two tails determine an arbitrary finite measure
modulo a constant multiple of Lebesgue measure, and they determine an
atomic measure uniquely.
\end{theorem}
\begin{proof}
The coefficient formula \eqref{eq:c-coeff} and the nonvanishing of
$B_{2j}$ show that a zero base tail implies
\[
 M_{m+2j}(\mu)=0,\qquad j\ge1.
\]
Independently of any reflection condition, the coefficient in the first
jet at $t^{m+2j}$ is a nonzero scalar multiple of
$M_{m+2j+1}(\mu)$, because \eqref{eq:trivial-zero-derivative} is nonzero.
Thus a zero first-jet tail also implies
$M_{m+2j+1}(\mu)=0$. Together, these say $M_n(\mu)=0$ for every
$n\ge m+2$. Eventual termination gives the same statement for all
sufficiently large $n$.

Let $\mathcal E(z)=\int\e^{az}\dd\mu(a)$. This is entire, since
$0<a<1$ and $\mu$ is finite. The Bernoulli generating function gives
\[
 \frac{z\mathcal E(z)}{\e^z-1}
       =\sum_{n\ge0}\frac{M_n(\mu)}{n!}z^n
       =P(z),\qquad |z|<2\pi,
\]
where $P$ is a polynomial. Multiplying through and applying the
identity theorem yields
\begin{equation}
 \mathcal E(z)=\frac{\e^z-1}{z}P(z)\quad\text{for all }z\in\C,
 \label{eq:measure-entire}
\end{equation}
with the removable value at $z=0$ understood. At every nonzero integer
$n$, this gives $\muh(n)=\mathcal E(2\pi\ii n)=0$.
The zeroth coefficient is $\alpha=\mu((0,1))$.

Regarding $(0,1)$ as the circle with its endpoint removed, $\mu$ and
$\alpha\lambda$ therefore agree on all trigonometric polynomials.
Those polynomials are uniformly dense in the continuous functions on
the circle. Finite measures that agree on them agree on all continuous
functions and hence are identical. Thus $\mu=\alpha\lambda$.
Conversely, integrating \eqref{eq:Bernoulli} over $a\in(0,1)$ shows
$M_n(\lambda)=0$ for all $n\ge1$, so all the high tails vanish for
$\lambda$. Equation \eqref{eq:Haar} gives its analytic interpretation.
A nonzero finite atomic measure cannot be a multiple of the nonatomic
measure $\lambda$, proving the last assertion.
\end{proof}

This result is linear injectivity modulo a completely identified
one-dimensional component. It does not depend on assuming the base
tail is already zero. For a measure in the reflection nullspace,
Theorem~\ref{thm:activation} gives the stronger interpretation of the
first-jet data as the actual hidden exponential function.

\begin{proposition}[General-measure reflection classification]
\label{prop:measure-reflection}
Let $\alpha=\mu((0,1))$. The base tail at order $m$ terminates if and
only if
\begin{equation}
 R_*(\mu-\alpha\lambda)=\eps_m(\mu-\alpha\lambda).
 \label{eq:measure-reflection}
\end{equation}
\end{proposition}
\begin{proof}
Integrate the normally convergent Borel kernel
\eqref{eq:base-kernel}. As in Theorem~\ref{thm:atomic}, termination
forces every residue to vanish. M\"obius inversion then makes all
positive Fourier coefficients of
$\nu=\mu+(-1)^mR_*\mu$ vanish. Its reflection parity supplies the
negative coefficients as well. By Fourier uniqueness,
$\nu=(1+(-1)^m)\alpha\lambda$. Rearranging gives
\eqref{eq:measure-reflection}. Conversely, Lebesgue measure contributes
no tail and the remaining reflection component cancels the Bernoulli
coefficients term by term.
\end{proof}

Theorem~\ref{thm:one-jet} is therefore sharp in two different senses:
no nonzero atomic measure is invisible to the pair of tails, but a
Lebesgue component is invisible to every order jet. The explicit
integral \eqref{eq:Haar} prevents an unjustified extension of atomic
injectivity to all finite measures without this quotient.

\section{Low-order identities and diagnostic examples}
\label{sec:examples}

\subsection{Order one: a Gamma reflection constant}
Using $\zeta'(0,a)=\log\Gamma(a)-\tfrac12\log(2\pi)$ and the Gamma
reflection formula, Theorem~\ref{thm:reflection} gives
\begin{equation}
 \begin{split}
 F_1(t,a)+F_1(t,1-a)
 &=\frac{\pi^2}{3t}-\log(2\sin\pi a)-\frac t2B_2(a)\\
 &\quad+2\sum_{d\ge1}\frac{\cos(2\pi da)}{d(\e^{4\pi^2d/t}-1)}.
 \end{split}
 \label{eq:m1}
\end{equation}
This is in the established $q$-product completion family, not an
independent priority claim. Its fixed-parity order derivative has the
finite part
\begin{equation}
 P_1^{[1]}(t,a)=\frac{2\zeta'(2)}t
  -(\log t+\gamma)\log(2\sin\pi a)
  -\frac{\zeta''(0,a)+\zeta''(0,1-a)}2
  +tB_2(a)\zeta'(0).
 \label{eq:m1-jet-P}
\end{equation}
Here $\zeta'(0)=-\tfrac12\log(2\pi)$. The convergent extra block is
\begin{equation}
 \Xi_1(t,a)=2\sum_{d\ge1}
 \frac{\cos(2\pi da)\log(t/(2\pi d))}{d(\e^{4\pi^2d/t}-1)}.
 \label{eq:m1-Xi}
\end{equation}
Equations \eqref{eq:m1-jet-P}, \eqref{eq:m1-Xi}, and
\eqref{eq:median-I} constitute an exact identity for a logarithmically
weighted Lambert sum involving Hurwitz-zeta second derivatives.

\subsection{A nonzero flat block with a zero base Borel kernel}
At $a=1/2$, \eqref{eq:m1} becomes
\begin{equation}
 2F_1(t,1/2)=\frac{\pi^2}{3t}-\log2+\frac{t}{24}
              +2\sum_{d\ge1}\frac{(-1)^d}{d(\e^{4\pi^2d/t}-1)}.
 \label{eq:half}
\end{equation}
The normalized base Borel kernel is identically zero. Nevertheless,
the last sum begins with $-2Q$, so it is not identically zero.
The first order derivative has double poles and its jump is exactly
$\pi\ii$ times that last sum. At $t=2$, the hidden remainder is of
order $10^{-9}$; the delivered verification checks it by subtracting
the finite core at 90-digit working precision, not by comparing only
power-series coefficients.

\subsection{Order two: an antisymmetric identity}
The next case has no $t^{-1}$ term:
\begin{equation}
 \begin{split}
 F_2(t,a)-F_2(t,1-a)
 &=\zeta(2)(1-2a)
       -t\bigl(\zeta'(-1,a)-\zeta'(-1,1-a)\bigr)
       +\frac{t^2}{6}B_3(a)\\
 &\quad-\frac t\pi\sum_{d\ge1}
                     \frac{\sin(2\pi da)}{d^2(\e^{4\pi^2d/t}-1)}.
 \end{split}
 \label{eq:m2}
\end{equation}
The order-one and order-two formulas exhibit the parity switch:
cosines survive in the hidden block for odd $m$, and sines survive
for even $m$. At $m$ even and $a=1/2$, the reflected observable is
identically zero. This is a zero measure after antisymmetric pairing,
not an exception to atomic detection.

\subsection{Direct coefficient reconstruction}
Suppose an admissible atomic combination has hidden series
$\sum b_NQ^N$. The first recovered Fourier coefficients are
\begin{align*}
 \muh(1)&=b_1,\qquad
 \muh(2)=2^m(b_2-b_1),\qquad
 \muh(3)=3^m(b_3-b_1),\\
 \muh(4)&=4^m(b_4-b_2),\qquad
 \muh(6)=6^m(b_6-b_3-b_2+b_1).
\end{align*}
These exact identities separate the divisor convolution from the
geometric Fourier problem. For a midpoint grid with $r=3$, the first
two recovered modes vanish and $\muh(3)=-3$; the first action is
$12\pi^2$, not $4\pi^2$.

\section{Source audit and corrections that matter}
\label{sec:audit}

\subsection{A confirmed cross-reference mismatch}
At the pinned commit, the source of \cite{repoq}, near its labels
\texttt{eq:FR60a} and \texttt{eq:FR60b}, describes the equations of
Fantini--Rella Conjecture~5 as ``(60a)--(60b)'' in arXiv:2506.08265v2.
The inspected v2 PDF, dated 1 April 2026, displays these as
\emph{(3.36a)--(3.36b)}, on its printed page 17 \cite{fr}.
The proposed correction is to change that external equation reference
to (3.36a)--(3.36b), retaining the repository's internal LaTeX labels
if desired. This is a version-specific documentary correction, not a
refutation of the mathematical identities. The delivery includes the
exact target sentence and replacement in a separate note.

\subsection{Double poles require a different finite part}
A potentially consequential extension error would be to apply ordinary
Cauchy principal-value integration directly to
\eqref{eq:jet-kernel}. For a nonzero local coefficient $c$, the symmetric
integral of $c/(\xi-A)^2$ outside an $\epsilon$-neighborhood diverges
like $2c/\epsilon$. The Laplace factor does not remove this divergence.
The repair is either a Hadamard finite part or the integrated-by-parts
simple-pole kernel of Lemma~\ref{lem:I}.

This is a warning for the new order-jet problem. The inspected
$q$-Gamma report \cite{repoq} uses simple poles and already distinguishes
its symmetric kernel from its exponential completion. It is not being
accused of making the double-pole error.

\subsection{Specialization must not replace analytic differentiation}
Theorem~\ref{thm:reflection} holds at integer $m$. Differentiating its
right side as though $m$ were a free complex variable is unjustified:
there is no specified holomorphic family of the parity sign
$\eps_m$, and the finite part changes at resonance. The correct
procedure is to differentiate $F_s+\eps_mF_s\circ R$ with the sign
fixed, on the initial Mellin contour, and only then specialize.
The missing tangent term in \eqref{eq:FE-jet-product} would be lost by
a naive differentiation of the specialized reflection formula.
That term is precisely the Borel-summed block of the first jet.

\subsection{Other distinctions retained in this report}
The zero base Borel kernel in \eqref{eq:half} coexists with a nonzero
flat remainder. This distinction is already explicit in the earlier
Lambert and repository literature; it is not a newly discovered error
there. Likewise, the linear median used here must not be substituted
for an algebra-compatible median in nonlinear products or inverse
functions.

The Fourier kernels in this report assume real shifts. For nonreal
$a$, a coefficient such as $\sin(2\pi da)$ can grow exponentially in
$d$, so the displayed Dirichlet expansions generally cease to converge.
One cannot claim complex-parameter uniformity by simply inserting
complex $a$ into them. Endpoint atoms at $a=0$ are also outside the
stated hypotheses. Finally, the Lebesgue counterexample
\eqref{eq:Haar} shows why the atomic injectivity theorem requires either
atomicity or the explicit quotient in Theorem~\ref{thm:one-jet}.

\section{Reproducibility and verification scope}
\label{sec:verification}

The mathematical arguments above prove infinite identities by contour
shifts, functional equations, normal convergence, and finite-dimensional
uniqueness. The accompanying program performs a different task: it
checks coefficients, constants, orientations, and representative
numerical identities that are especially vulnerable to transcription
errors. It is not a proof assistant and it does not establish priority.

\subsection{Exact symbolic and arithmetic checks}
The delivered run records \textbf{732 passing exact assertions} in
SymPy~1.14.0. They include Bernoulli reflection cancellation and the
activated coefficients for $1\le m,j\le8$; rational-kernel residues and
double-pole principal parts; the $\pi\ii$ normalization through
$m=24$; generic collision jets through derivative order six;
Fourier--M\"obius calculations in exact cyclotomic fields of moduli
$3$ through $10$; midpoint-grid delay; and the Bernoulli multiplication
formula. The free-variable symbolic tests are distinguished in the
record from finite arithmetic instances. These computations support,
but do not replace, the general proofs.

\subsection{High-precision diagnostics}
There are \textbf{61 passing numerical assertions}, primarily at
90 decimal digits using mpmath~1.3.0. The independent contour tests use
55 digits. The error recorded is
\begin{equation}
 \frac{|\mathrm{LHS}-\mathrm{RHS}|}
           {\max(1,|\mathrm{LHS}|,|\mathrm{RHS}|)}.
 \label{eq:error-metric}
\end{equation}
It is a scaled absolute floating-point error, not an interval enclosure.
The largest observed errors in the delivered run are summarized below.

\begin{center}
\small
\begin{tabular}{@{}p{75mm}rr@{}}
\toprule
Diagnostic & Count & Largest observed error\\
\midrule
Exact reflection completion & 36 & $2.46\times10^{-91}$\\
Exact first-order-jet completion & 10 & $5.84\times10^{-47}$\\
Acceleration refinement & 10 & $5.84\times10^{-47}$\\
Independent rotated-contour jump & 2 & rounded to zero\\
Independent linear-median kernel & 2 & $5.74\times10^{-57}$\\
Nonzero hidden remainder at $a=1/2$ & 1 & $2.55\times10^{-91}$\\
\bottomrule
\end{tabular}
\end{center}
The acceptance thresholds were $10^{-75}$ for reflection and the hidden
remainder, $10^{-38}$ for the full first-jet identities and refinement,
and $10^{-42}$ for the contour tests. A rounded zero in one diagnostic
is not a claim of exact numerical equality.

The reflection tests use $m=1,\ldots,6$, shifts
$a=1/3,1/2,\sqrt2/3$, and $t=2,5$. The first-jet tests use
$m=1,\ldots,5$, shifts $1/3,\sqrt2/3$, and $t=3+m/2$.
The original jet is computed from its convergent series with the exact
weight $-\log k$, not by differencing nearby orders.
The contour tests integrate the rational double-pole kernel on rays of
angles $\pm0.32$ radians at $\lambda=4,9$, and compare to
\eqref{eq:double-Laplace}. They check the sign convention independently
of the arithmetic summations.

\subsection{Why the median evaluation is accelerated}
Directly summing \eqref{eq:median-I} gives only an $N^{-3}$ tail.
For accurate evaluation, subtract finitely many terms of the asymptotic
expansion of $I$ from each summand and sum those terms globally using
\eqref{eq:d-coeff}. For any fixed integer $J\ge1$, this gives the
\emph{exact infinite-sum identity}
\begin{equation}
 \begin{split}
 t^m\Med\widetilde J_m
 &=t^m\sum_{j=1}^Jd_{m,j}(a)t^{2j}\\
 &\quad+\left(\frac{\ii t}{2\pi}\right)^{m-1}
 \sum_{N\ge1}E_m(N,a)
 \left(I(A_N/t)-\sum_{j=1}^J\frac{(2j)!}{(A_N/t)^{2j+1}}\right).
 \end{split}
 \label{eq:accelerated}
\end{equation}
Only the outer $N$-sum is numerically truncated. The delivered run uses
$J=12$, $N\le72$ and compares with $J=13$, $N\le80$.
This is finite subtraction, not an assertion that the factorial
asymptotic series converges. The numerical truncation and floating-point
rounding are not interval-certified; their agreement is a diagnostic.

The general-measure kernel theorem is established by its generating-
function proof. The numerical suite does not purport to test all finite
measures, unknown support recovery, or complex-shift analytic continuation.

\section{Further research questions}
\label{sec:future}

\subsection{Higher order jets and logarithmic Borel singularities}
Corollary~\ref{cor:collision-reflection} determines every collision
polynomial, but this report gives a complete Borel kernel and exact
completion only for the first order derivative. Higher derivatives of
$\zeta(s+w)$ introduce polygamma terms, higher derivatives of the
functional-equation trigonometric factor, and logarithmic divisor
weights. An explicit all-jet completion should be sought with a clear
separation of local collision polynomials, lateral discontinuities,
and independently identified flat blocks. The first task is to obtain
and prove the second-order kernel, rather than infer it by formal
specialization.

\subsection{Two shifted lattice directions}
Replacing the $k^{-s}$ factor by a second shifted lattice sum leads to
Mellin products of the form
$\Gamma(w)\zeta(w,a)\zeta(s+w,b)$. The two Hurwitz functional equations
create a matrix of reflection sectors rather than a single parity.
Which sectors terminate, and how many order jets are needed to detect
all finite atomic measures on the two-torus? The one-dimensional
Lebesgue blind component suggests that continuous product components
will have to be identified explicitly.

\subsection{Unknown shifts and inverse arithmetic data}
Theorem~\ref{thm:action-detection} gives a finite reconstruction on known
nodes. Unknown-node recovery suggests a Prony-type reconstruction from
Fourier coefficients obtained by M\"obius inversion. Sharp sample
counts, treatment of collisions and multiplicities, and certified
recovery from approximate Stokes data remain separate problems. Exact
identifiability should be settled before numerical stability is claimed.

\subsection{Complex shifts and moving endpoints}
A kernel holomorphic in a complex neighborhood of $a$ must replace the
real Fourier representation; a local meromorphic or integral
representation is a plausible route. A second problem is the
nonuniform limit in which $a$ approaches $0$ or $1$ together with $t$.
The fixed-shift Gamma core then changes scale. Neither extension follows
from the normal-convergence statements used here.

\subsection{Nonlinear inversion after completion}
The exact functions $N_m$ and $\Xi_m$ can be transported through a
regular inverse or a critical inverse chart. Existing repository work
already treats the distinction between formal inversion and inversion
of the physically completed function \cite{repoq}. A useful new
application would choose an actual critical point of a Lambert--
polylogarithm combination and compute its first nonzero ramified
exponential sectors from the completed forward function. Merely applying
a formal inverse theorem to an unspecified flat completion would not
settle that problem.

\subsection{Formalization and reusable exact certificates}
The Bernoulli reflection, divisor convolution, M\"obius inversion,
Vandermonde uniqueness, and midpoint-grid distribution are finite
algebraic targets for formalization. They should be separated from the
analytic Mellin shifts and finite-part integrals. A later analytic
formalization would need parameter differentiation under normal
convergence, a contour-shift theorem with stated growth bounds, and a
precise double-pole finite-part convention. The supplied Python checks
are not substitutes for those proof-assistant declarations.

\section*{Conclusion}
\addcontentsline{toc}{section}{Conclusion}
Reflection can remove an entire perturbative tail without removing the
analytic function's exponentially small remainder. For shifted
Lambert--polylogarithm sums, the first order derivative restores exactly
the complementary Fourier information, now at double Borel poles.
Mellin inversion identifies the derivative's missing logarithmic dual
block and proves the factor-$\pi\ii$ recovery law. Finite atomic
reflection rigidity turns that analytic statement into a finite
identification theorem, while the general-measure result isolates
Lebesgue measure as the sole simultaneous blind component.

\appendix
\section{A compact dependency and status ledger}
\label{app:ledger}

\begin{center}
\small
\begin{tabular}{@{}p{45mm}p{58mm}p{35mm}@{}}
\toprule
Result & Main dependencies & Evidence status\\
\midrule
All-order collision formula & Mellin differentiation and Laurent coefficients
 & Proof; finite jet checks\\[3pt]
Simple-pole Borel kernel & Bernoulli Fourier expansion and a geometric series
 & Proof; exact normalization checks\\[3pt]
Exact reflected duality & Hurwitz/Riemann functional equations and Mellin inversion
 & Proof; 36 numerical instances\\[3pt]
Double-pole activation & Trivial-zero derivatives and contour orientation
 & Proof; independent contour tests\\[3pt]
Exact first-jet completion & Mellin integration by parts and the $I$ kernel
 & Proof; 10 numerical instances\\[3pt]
Atomic reflection and action bound & M\"obius inversion and Vandermonde uniqueness
 & Proof; exact cyclotomic tests\\[3pt]
General one-jet kernel & Bernoulli generating function and Fourier uniqueness
 & Proof; not numerically certified\\
\bottomrule
\end{tabular}
\end{center}

No theorem in this ledger is asserted to be Lean-formalized. The
literature comparison, source cross-reference check, symbolic tests,
numerical diagnostics, and PDF build are distinct evidence scopes.
The source report's broader claims about other conjectures are not
required by any proof in this article.

\section{Rebuilding the package}
The top-level \texttt{Makefile} builds this article with three ordinary
pdfLaTeX passes and runs the exact and numerical checks separately.
The verification script accepts \texttt{--mode exact},
\texttt{--mode numerical}, or \texttt{--mode all}, and writes only the
explicit output path. Its exact dependencies are listed in
\texttt{requirements.txt}. A typical replay is
\begin{verbatim}
python verification/verify.py --mode exact \
  --output verification/exact_results.json
python verification/verify.py --mode numerical \
  --output verification/numerical_results.json
make pdf
\end{verbatim}
The delivered records were produced with Python~3.13.5, SymPy~1.14.0,
and mpmath~1.3.0. Replays should be made on a copy of the package if the
original recorded outputs are to remain unchanged. The report was
prepared with ChatGPT assistance; no independent referee review is
implied by the internal proof-review notes.

\begin{thebibliography}{9}
\addcontentsline{toc}{section}{References}

\bibitem{dlmf}
NIST Digital Library of Mathematical Functions.
Sections 24.8 (Bernoulli Fourier series), 25.4 (Riemann-zeta functional
equation), 25.11 (Hurwitz zeta), and 25.12 (polylogarithms).
\url{https://dlmf.nist.gov/24.8},
\url{https://dlmf.nist.gov/25.4},
\url{https://dlmf.nist.gov/25.11},
\url{https://dlmf.nist.gov/25.12}.
Consulted for the classical identities used in this report.

\bibitem{dk}
D.~Dorigoni and A.~Kleinschmidt,
\emph{Resurgent expansion of Lambert series and iterated Eisenstein
integrals}, 2020.
\url{https://arxiv.org/abs/2001.11035}.
In particular, the scalar odd-integer degeneration and surviving
nonperturbative sectors are discussed in Section~2.2.

\bibitem{fr}
V.~Fantini and C.~Rella,
\emph{Modular resurgence, $q$-Pochhammer symbols, and quantum operators
from mirror curves}, arXiv:2506.08265v2, 1 April 2026.
\url{https://arxiv.org/abs/2506.08265v2}.
The inspected Conjecture~5 appears on printed page~17,
Eqs.~(3.36a)--(3.36b).

\bibitem{repoq}
ProveIt research report,
\emph{Critical Stokes Reversion and $q$-Gamma Cores}, source dated
4 October 2026, pinned to commit \pin.
Repository path:
\path{Analysis/Transseries/docs/series-and-transseries/}%
\path{Mellin_Borel_Completion_q_Gamma_Critical_Reversion/}%
\path{critical_stokes_q_reversion.tex}.
The inspected ranges include the title/scope, simple-pole kernel,
Mellin completion proof, and rational-weight criterion.
Repository: \url{https://github.com/VladimirReshetnikov/ProveIt}.

\bibitem{repoindex}
V.~Reshetnikov and ProveIt contributors,
\emph{ProveIt repository documentation}.
Inspected paths:
\path{Analysis/Polylogarithms/README.md},
\path{Analysis/Transseries/README.md},
and \path{docs/incoming/README.md}, together with the incoming directory
listing. Live indexes were read on 12 October 2026; the detailed
mathematical comparison is pinned separately in \cite{repoq}.
\url{https://github.com/VladimirReshetnikov/ProveIt}.

\end{thebibliography}
\end{document}
